\section{The gluing problem in categorical generality}
\label{sec:gluing:abstract}
\label{ch:gluing-part-i}
\subsection{The abstract set-up}
\label{sec:glue-abstract-setup}

Fix a Calabi--Yau variety $X$ of complex dimension $d$, smooth or
with at worst Gorenstein canonical singularities, equipped with a
nowhere-vanishing holomorphic volume form $\Omega_X \in
H^{d,0}(X)$. The cohomological Hall algebra $\CoHA(X)$ of
Kontsevich--Soibelman arXiv:1006.2706 is an associative graded
algebra built from vanishing-cycle cohomology of moduli of objects
in a Calabi--Yau category attached to $X$. The construction is
\emph{local-to-global by design}: critical cohomology lives on a
stack glued from local charts, and the Hall convolution is defined
in terms of local extension correspondences.

The local input is a quiver with potential $(Q_\alpha, W_\alpha)$
presenting a Calabi--Yau $A_\infty$-algebra on a chart
$U_\alpha \subset X$; the global output is a sheaf of algebras whose
sections over $U_\alpha$ recover the local critical CoHA. The
passage is governed by the combinatorial datum it must resolve: the
\emph{Derksen--Weyman--Zelevinsky mutation groupoid}
$\mathrm{Mut}(X)$, whose objects are quiver-with-potential
presentations of $X$, whose 1-morphisms are mutations
(arXiv:0704.0649, Theorem~1.7), and whose 2-morphisms are
intertwiners of the induced Keller--Yang Fourier--Mukai kernels
(arXiv:0906.0761, Theorem~3.2).

\emph{Truncation level of $\mathrm{Mut}(X)$.} Presentations are
sets (0-truncated); mutations between presentations form a
non-trivial groupoid of Fourier--Mukai kernels (1-truncated
non-trivially, since distinct mutations produce non-equal kernels);
intertwiners between such kernels are \emph{contractible} by the
Keller--Yang canonicity theorem (arXiv:0906.0761, Proposition~3.3:
the silting-object lift is unique up to unique isomorphism). The
groupoid of intertwiners-of-intertwiners is therefore contractible,
and $\mathrm{Mut}(X)$ is a $(2,1)$-truncated simplicial set in the
sense of Lurie, HTT~\S2.3.4.

\emph{Three consequences of $(2,1)$-truncation.} Locality: vanishing
cycles are étale-local on the target of $W$, so the 0-cells of
$\mathrm{Mut}(X)$ supply local CoHA data. Covariance: every
mutation $(Q_\alpha, W_\alpha) \to (Q_\beta, W_\beta)$ is a 1-cell
in $\mathrm{Mut}(X)$ lifting to a canonical Fourier--Mukai kernel
(Keller--Yang~\S3.2), so the overlap data live at the 1-cell level.
Derived invariance: Keller--Yang canonicity is a 2-cell
contractibility statement, recording that the gluing result is
insensitive to the mutation path chosen.

The $(2,1)$-truncation of $\mathrm{Mut}(X)$ is therefore the
minimal truncation that captures all three data: higher truncations
record nothing further (by Keller--Yang canonicity), lower
truncations lose the Fourier--Mukai kernel data (by the
non-triviality of mutation at the 1-cell level). The abstract
object constructed below is the value on $\mathrm{Mut}(X)$ of a
sheaf in the Toën--Vezzosi derived-descent sense
(arXiv:math/0404373, Section~1.3).

\begin{definition}[Critical CoHA on a smooth chart, after Kontsevich--Soibelman 2011]
\label{def:critical-coha-local}
\ClaimStatusDefinition
Let $Y$ be a smooth quasi-projective variety with an action of a
reductive group $G$ and let $f: Y \to \bC$ be a $G$-invariant
regular function. The \emph{critical cohomological Hall algebra}
of $(Y, f, G)$ is the graded vector space
\[
 \sH(Y, f, G) \;=\; H^{*}_{G}\!\bigl(Y,\, \varphi_f\bigr)
\]
where $\varphi_f$ is the sheaf of vanishing cycles of $f$; the
convolution product, when $(Y, f, G)$ is equipped with a choice of
extension correspondence, is
$m(\alpha \otimes \beta) = q_* p^*(\alpha \boxtimes \beta)$ along
\[
 Y' \times Y''
 \;\xleftarrow{p}\;
 Y'_{\subset}
 \;\xrightarrow{q}\;
 Y
\]
twisted by the Joyce--Song sign
(Kontsevich--Soibelman arXiv:1006.2706, Definition~2; Davison
arXiv:1311.7172, Theorem~A for regularity of $\varphi_f$ under the
equivariant structure; Davison--Meinhardt arXiv:1601.02479,
Theorem~A for associativity and the dimensional-reduction
isomorphism).

The \emph{quiver-with-potential form} is the specialisation
$(Y, f, G) = (\Rep(Q, \mathbf d), W_{\mathbf d}, G_{\mathbf d})$
with $G_{\mathbf d} = \prod_i GL(d_i)$, $W_{\mathbf d}$ the trace
of the potential, and the extension correspondence given by
short-exact-sequence stacks $0 \to V' \to V \to V'' \to 0$ in
$\Rep(Q)$. When this specialisation is available, $\sH(Q, W) =
\bigoplus_{\mathbf d} \sH(\Rep(Q, \mathbf d), W_{\mathbf d}, G_{\mathbf d})$
is $\bN^{Q^0}$-graded.
\end{definition}

\begin{proposition}[When a chart admits quiver presentation]
\label{prop:chart-quiver-presentation}
\ClaimStatusProposition
Let $U \subset X$ be an affine open in a Calabi--Yau threefold.
The critical CoHA $\sH(U) = \sH(Y_U, f_U, G_U)$ of
Definition~\ref{def:critical-coha-local} admits a
quiver-with-potential presentation, i.e.\ an isomorphism $\sH(U)
\cong \sH(Q, W)$ for some $(Q, W)$, under either of the following
\emph{global} hypotheses on $U$:
\begin{enumerate}
\item[(VdB)] $U$ admits a \emph{Van den Bergh tilting bundle}: a
 coherent sheaf $T$ on $U$ with $\mathrm{Ext}^{>0}(T, T) = 0$,
 $\mathrm{End}(T)$ Noetherian, and the functor
 $R\Hom(T, -): \Db(\Coh U) \to \Db(\mathrm{End}(T)\text{-mod})$ an
 equivalence (Van den Bergh arXiv:math/0211064, Theorem~A). Then
 $(Q, W)$ is the Ginzburg quiver of $\mathrm{End}(T)$, and
 $U = \Spec \Jac(Q, W)$ (Ginzburg arXiv:math/0612139,
 Theorem~3.6.4).
\item[(DM)] $U$ is toric without compact divisors, and the
 Rapcak--Soibelman--Yang--Zhao presentation of
 Davison--Meinhardt arXiv:1601.02479 Theorem~A applies. Then
 $(Q, W)$ is the Jacobi quiver of the toric fan.
\end{enumerate}
Neither hypothesis is automatic: (VdB) fails on compact
Calabi--Yau threefolds (a tilting bundle would force $H^0(U, \sO)$
to be Noetherian, incompatible with the Čech assembly of a
projective CY\textsubscript{3}); (DM) fails whenever $U$ carries
a compact 4-cycle.
\end{proposition}

\begin{proof}[Remark on the proof]
Both implications (VdB) $\Rightarrow$ Jacobi and (DM) $\Rightarrow$
toric Jacobi are by direct quotation of the cited theorems; the
negative statements are obstructions: any tilting bundle is
concentrated in degree 0, and toric compact-divisor support forces
the fan to have a compact cone, which then supports a non-zero
higher-direct-image obstructing tilting.
\end{proof}

\begin{remark}[Chart cover hypothesis, named explicitly]
\label{rem:scope-critical-coha-local}
The rest of the chapter fixes an open cover $\mathfrak{U} =
\{U_\alpha\}$ of $X$ with each $U_\alpha$ satisfying either (VdB)
or (DM). Such a cover exists when $X$ is either toric (DM, fine
enough fan refinement) or admits a local-to-global Van den Bergh
resolution on affine charts (Buchweitz--Hille arXiv:1404.3007 for
the flop-local case; still open in general for compact CY\textsubscript{3}).
The existence hypothesis is the chart-level input to the
gluing; it does not imply a global tilting bundle on $X$, which
would contradict compactness.
\end{remark}

\subsection{What ``sheaf of CoHAs'' means}
\label{sec:sheaf-of-cohas}

The sheaf-theoretic content of the critical CoHA is visible on each
individual chart: vanishing cycles are étale-local on the target of
$W$, and Borel--Moore homology extends along open inclusions. A
single chart $U$ therefore carries a locally-constant sheaf $\sH_U$
of graded algebras whose sections over an affine open $V \subset U$
with Jacobi presentation $\Jac(Q_V, W_V)$ recover $\sH(Q_V, W_V)$.
The question that organises this section is the following.

\emph{Question.} When two charts $U_\alpha$, $U_\beta$ overlap on
$U_{\alpha\beta} = U_\alpha \cap U_\beta$, each chart carries its
own presentation $(Q_\alpha, W_\alpha)$, $(Q_\beta, W_\beta)$, and
the overlap receives two \emph{a priori} distinct presentations.
Which canonical datum identifies them?

The answer is forced by three facts. First, on a Calabi--Yau
threefold any two quiver-with-potential presentations of the same
smooth affine Calabi--Yau threefold are related by a finite sequence
of Derksen--Weyman--Zelevinsky mutations
(arXiv:0704.0649, Theorem~1.7); each mutation induces a derived
equivalence of Jacobi categories (Keller--Yang arXiv:0906.0761,
Theorem~3.2). Second, each such derived equivalence preserves the
critical locus of $W$ up to an explicit Fourier--Mukai kernel on
the moduli stack. Third, Davison--Meinhardt arXiv:1601.02479
Theorem~A shows that this Fourier--Mukai transport sends the
critical CoHA of one presentation to the critical CoHA of the other
by an \emph{algebra isomorphism} determined canonically by the
tilting object. The overlap datum is thus a \emph{tilting cocycle}
in Fourier--Mukai kernels.

Let $\mathrm{Alg}^{\mathrm{FM}}_{(2,1)}$ denote the
$(2,1)$-category whose objects are $\bN^\bullet$-graded associative
algebras of the form $\sH(Q, W)$, whose 1-morphisms are invertible
Fourier--Mukai transports $\Phi_K: \sH(Q_1, W_1) \xrightarrow{\sim}
\sH(Q_2, W_2)$ induced by derived-equivalence kernels $K$ between
the Jacobi categories, and whose 2-morphisms are kernel intertwiners
modulo Keller--Yang contractibility (arXiv:0906.0761,
Proposition~3.3). This is well-defined by Keller--Yang plus the
$(\infty,1)$-nerve construction Lurie HTT~\S1.2.6.1.

\begin{definition}[Simplicial $(2,1)$-sheaf of CoHAs]
\label{def:simplicial-sheaf-cohas}
\ClaimStatusDefinition
Let $\mathfrak{U} = \{U_\alpha\}_{\alpha \in A}$ be an analytic
open cover of $X$ with each $U_\alpha$ satisfying (VdB) or (DM) of
Proposition~\ref{prop:chart-quiver-presentation} and carrying a
chosen presentation $(Q_\alpha, W_\alpha)$. Write
$\Cech^\bullet(\mathfrak{U})$ for the nerve:
$\Cech^n(\mathfrak{U}) = \coprod_{\alpha_0 < \cdots < \alpha_n}
U_{\alpha_0 \cdots \alpha_n}$. A \emph{simplicial $(2,1)$-sheaf of
CoHAs} on $(X, \mathfrak{U})$ is a homotopy-coherent simplicial
diagram
\[
 \sH_\bullet: \Delta^{\mathrm{op}} \longrightarrow
 \mathrm{Alg}^{\mathrm{FM}}_{(2,1)}
\]
in the sense of Lurie HTT~\S1.2.6.1 (homotopy-coherent nerve),
subject to the following simplex-by-simplex data:
\begin{enumerate}
\item[(S1)] \emph{0-simplices.} For each $\alpha$,
 $\sH_0|_{U_\alpha} = \sH(Q_\alpha, W_\alpha)$ per
 Definition~\ref{def:critical-coha-local}.
\item[(S2)] \emph{1-simplices.} For each pair $\alpha < \beta$,
 the face 1-morphisms
 $\partial_0, \partial_1: \sH_1|_{U_{\alpha\beta}}
 \rightrightarrows \sH_0|_{U_\alpha},\, \sH_0|_{U_\beta}$ are
 invertible Fourier--Mukai transports $\Phi_{K_{\alpha\beta}}$
 induced by DWZ mutation kernels (arXiv:0704.0649, Theorem~1.7;
 Keller--Yang arXiv:0906.0761, Theorem~3.2).
\item[(S3)] \emph{2-simplices (cocycle).} For each triple
 $\alpha < \beta < \gamma$, the 2-simplex
 $\sH_2|_{U_{\alpha\beta\gamma}}$ carries an invertible 2-morphism
 \[
 \eta_{\alpha\beta\gamma}:
 \Phi_{K_{\beta\gamma}} \circ \Phi_{K_{\alpha\beta}}
 \;\xrightarrow{\;\sim\;}\;
 \Phi_{K_{\alpha\gamma}}
 \]
 in $\mathrm{Alg}^{\mathrm{FM}}_{(2,1)}$, witnessing
 $K_{\beta\gamma} \ast K_{\alpha\beta} \simeq K_{\alpha\gamma}$ up
 to the Keller--Yang canonical 2-isomorphism.
\item[(S4)] \emph{Contractible filling at $n \geq 3$.} Inner horns
 $\Lambda^n_k \hookrightarrow \Delta^n$ for $n \geq 3$ admit unique
 fillers: by Keller--Yang Proposition~3.3, 2-morphisms in
 $\mathrm{Alg}^{\mathrm{FM}}_{(2,1)}$ are isomorphisms and
 isomorphisms-of-isomorphisms are unique, so no further choice is
 required at degree $\geq 3$.
\end{enumerate}
\end{definition}

\begin{remark}[Strictification obstructed; $(2,1)$-truncation optimal]
\label{rem:strictness-fails}
A strict simplicial diagram demanding equalities
$\partial_j \partial_i = \partial_i \partial_{j+1}$ at the
1-morphism level is unattainable: the DWZ composite
$K_{\beta\gamma} \ast K_{\alpha\beta}$ agrees with $K_{\alpha\gamma}$
only up to the Keller--Yang canonical 2-isomorphism, and forcing
equality would require a splitting of the Keller--Yang contractible
space of intertwiners (no canonical section exists). Item~(S3)
upgrades the equality to an explicit invertible 2-morphism
$\eta_{\alpha\beta\gamma}$. Truncation at $(2,1)$-level is then
forced by (S4): no $(n,1)$-truncation with $n \geq 3$ records
anything beyond what Keller--Yang already determines uniquely
(Lurie HTT~Proposition~2.3.4.18 on the uniqueness of truncations;
Toën--Vezzosi arXiv:math/0404373, Section~1.3 on $(2,1)$-stacks as
the home for derived descent of algebra-valued presheaves).
\end{remark}

\begin{proposition}[Descent reduction via homotopy totalisation]
\label{prop:descent-reduction}
\ClaimStatusProposition
Given a simplicial $(2,1)$-sheaf of CoHAs $\sH_\bullet$ satisfying
(S1)--(S4) of Definition~\ref{def:simplicial-sheaf-cohas}, the
$(2,1)$-categorical \emph{homotopy totalisation}
\[
 \sH(X) \;=\; \mathrm{Tot}\!\bigl(\sH_\bullet\bigr)
 \;=\;
 \mathrm{holim}_{[n] \in \Delta}\, \sH_n
 \;\simeq\;
 \int_{[n] \in \Delta} \sH_n^{\Delta^n}
\]
(the right-hand expression is the end computing the cosimplicial
totalisation; Lurie HTT~\S4.2.4, formula (4.2.4.1); equivalent by
Lurie HTT~Corollary~4.2.4.8 to the homotopy limit over the
opposite simplex category) is a well-defined torsor-pointed
$\bN^\bullet$-graded associative algebra,
independent of the cover $\mathfrak{U}$ and of the choice of
presentations $(Q_\alpha, W_\alpha)$:
\begin{enumerate}
\item Cover independence is by cofinality of common refinements in
 the poset of open covers, combined with Toën--Vezzosi
 arXiv:math/0404373 Section~1.3 on homotopy-theoretic descent of
 presheaves valued in $(2,1)$-categories.
\item Presentation independence is by Davison--Meinhardt
 arXiv:1601.02479 Theorem~A: any two chart-wise presentations are
 related by a DWZ mutation cocycle, and the Fourier--Mukai
 transport along a mutation cocycle induces an invertible
 1-morphism in $\mathrm{Alg}^{\mathrm{FM}}_{(2,1)}$.
\item Canonicity of $\sH(X)$ after fixing the Keller--Yang/Tamarkin
coherent 2-isomorphism torsor is by the
 Francis--Gaitsgory chiral-Koszul totalisation arXiv:1103.5803,
 Theorem~1.2.4 (specialised from Ran-space totalisation to the
 nerve $\Cech^\bullet(\mathfrak{U})$, viewed as a finite
 subdiagram of the Ran space supported on the cover).
\end{enumerate}
The algebra $\sH(X)$ is the \emph{global critical CoHA} of $X$.
\end{proposition}

The load-bearing hypothesis is the existence of a simplicial
$(2,1)$-sheaf $\sH_\bullet$ satisfying (S1)--(S4) for $(X,
\mathfrak{U})$: the 2-simplex filler (S3) must be constructible
geometrically, not merely posited. Four mechanisms for constructing
it are known, each tied to an equivariance stratum of $X$; their
taxonomy occupies Section~\ref{sec:four-stratum-taxonomy}.

\subsection{The gluing 2-category \texorpdfstring{$\GlCoHA(X)$}{GluedCoHA(X)}}
\label{sec:glue-2cat}

The existence and uniqueness of $\sH(X)$ asserted in
Proposition~\ref{prop:descent-reduction} is cleanest to state in
the $(2,1)$-category of all possible sheaves-of-CoHAs on $X$.
Remark~\ref{rem:strictness-fails} recorded that the truncation of
the ambient category at $(2,1)$ is forced by the data; before the
definition of $\GlCoHA(X)$ itself, the same argument is re-run at
the level of the category of simplicial sheaves, to rule out the
alternative truncation levels $(1,1)$, $(2,2)$, and $(\infty,1)$
globally.

\begin{proposition}[$(2,1)$-truncation of $\GlCoHA(X)$ is forced]
\label{prop:glcoha-truncation-forced}
\ClaimStatusProposition
Let $X$ be a Calabi--Yau variety admitting a chart-wise
quiver-with-potential cover. Among the four candidate truncation
levels $(n, m) \in \{(1,1),\, (2,1),\, (2,2),\, (\infty, 1)\}$ for a
category $\GlCoHA(X)$ whose objects are simplicial sheaves of CoHAs
on $X$, exactly one, $(n, m) = (2, 1)$, records the full gluing
datum losslessly.
\begin{enumerate}
\item[(i)] $(1,1)$ loses the Keller--Yang tilting cocycle at the
 1-cell level: forgetting 2-morphisms identifies the
 distinct-but-isomorphic Fourier--Mukai kernels
 $K_{\beta\gamma} \ast K_{\alpha\beta}$ and $K_{\alpha\gamma}$ into
 a single connected component, at which point (S3) of
 Definition~\ref{def:simplicial-sheaf-cohas} collapses and the
 cocycle condition becomes strict. Strictness is obstructed
 (Remark~\ref{rem:strictness-fails}).
\item[(ii)] $(2,1)$ records the tilting cocycle as an
 invertible 1-morphism and records the Keller--Yang canonical
 isomorphism $\eta_{\alpha\beta\gamma}$ of (S3) as a contractible
 2-morphism. No higher data intrudes by (S4).
\item[(iii)] $(2,2)$ requires 2-morphisms to admit their own
 non-trivial automorphisms. The intertwiner space of any two
 Keller--Yang 2-isomorphisms is a $\mathrm{GRT}_1(\bQ)$-torsor
 over the Tamarkin associator (Kontsevich 1999
 \emph{Lett. Math. Phys.}~48, Theorem~3;
 Tamarkin arXiv:math/0411168, Theorem~A). The group
 $\mathrm{GRT}_1(\bQ)$ is pro-unipotent (Drinfeld, \emph{Algebra
 i Analiz} 2 (1990), \S 5, Theorem~A$'$; Furusho arXiv:0907.4715,
 Theorem~0.1), hence contractible as a pro-algebraic group; any
 $\mathrm{GRT}_1(\bQ)$-torsor is therefore contractible, so
 $\pi_k = 0$ for every $k \geq 1$. No non-trivial
 2-morphism-of-2-morphisms arises.
\item[(iv)] $(\infty, 1)$ coincides with $(2, 1)$ under canonical
 inclusion: the homotopy-coherent nerve of the simplicial diagram
 of Definition~\ref{def:simplicial-sheaf-cohas} is 2-truncated, so
 Lurie HTT~Proposition~2.3.4.18 identifies the $(\infty, 1)$-limit
 with the $(2, 1)$-limit.
\end{enumerate}
\end{proposition}

\begin{proof}
\emph{(i).} The obstruction of Remark~\ref{rem:strictness-fails} is
precisely the failure of strict cocyclicity; a $(1,1)$-category
forgets the 2-morphism $\eta_{\alpha\beta\gamma}$ and consequently
demands strict equality of $K_{\beta\gamma} \ast K_{\alpha\beta}$
and $K_{\alpha\gamma}$. No canonical splitting of the Keller--Yang
2-isomorphism torsor exists (Keller--Yang arXiv:0906.0761,
Proposition~3.3: intertwiners are canonical as a coherent torsor, with
no canonical section), so the $(1,1)$-category
collapses non-equivalent gluing data into a single class.

\emph{(ii).} The 1-morphism space in
$\mathrm{Alg}^{\mathrm{FM}}_{(2,1)}$ is by construction a groupoid,
and the 2-morphism space is the groupoid of kernel intertwiners
modulo Keller--Yang contractibility. The simplicial diagram of
Definition~\ref{def:simplicial-sheaf-cohas} uses precisely these
1- and 2-morphisms; (S3) encodes the Keller--Yang 2-isomorphism and
(S4) asserts contractibility of fillers at $n \geq 3$. This is
exactly the $(2,1)$-truncation.

\emph{(iii).} Fix two Keller--Yang 2-isomorphisms
$\eta, \eta': K_{\beta\gamma} \ast K_{\alpha\beta}
\xrightarrow{\sim} K_{\alpha\gamma}$. The space of intertwiners
$\eta \Rightarrow \eta'$ is identified via Kontsevich--Tamarkin
formality with the mapping space of $L_\infty$-morphisms between
two instances of the Hochschild complex of the Ginzburg dg-algebra
$\Jac(Q, W)$; on the single connected component supported by the
Tamarkin associator this space is a $\mathrm{GRT}_1(\bQ)$-torsor
(Kontsevich 1999, Section~4; Tamarkin op.~cit., Section~5).
$\mathrm{GRT}_1(\bQ)$ is pro-unipotent (Drinfeld, \emph{Algebra i
Analiz} 2 (1990), Theorem~A$'$; Furusho arXiv:0907.4715,
Theorem~0.1): its rational homotopy type is that of a pro-finite
product of affine spaces $\prod_{k \geq 3, \text{odd}} \bA^1_\bQ$.
A torsor over a contractible pro-algebraic group is contractible,
so $\pi_k = 0$ for $k \geq 1$: no non-trivial automorphism of a
2-morphism arises.

\emph{(iv).} The homotopy-coherent nerve of the $(2,1)$-simplicial
diagram inherits its truncation level from the target category:
since $\mathrm{Alg}^{\mathrm{FM}}_{(2,1)}$ is 2-truncated by
definition, so is its diagram category. Lurie HTT~\S2.3.4 classifies
2-truncated $(\infty, 1)$-categories as $(2, 1)$-categories, and
HTT~Proposition~2.3.4.18 asserts that homotopy limits computed in
either are canonically identified.
\end{proof}

\begin{remark}[Truncation is a property of the data, not a choice]
\label{rem:truncation-property-not-choice}
Proposition~\ref{prop:glcoha-truncation-forced} is not a choice of
convention: the truncation level of $\GlCoHA(X)$ is dictated by the
content of the gluing cocycle. Raising the truncation above
$(2,1)$ records no additional information; lowering it below
$(2,1)$ discards the Keller--Yang canonicity that makes the global
CoHA torsor-pointed and independent of presentations. An attempt to work at
$(3,1)$ or $(\infty,1)$ is equivalent to $(2,1)$ up to canonical
inclusion, paid for by additional formalism. An attempt to work at
$(1,1)$ is inadequate: it cannot distinguish ``gluing succeeded
canonically'' from ``gluing succeeded non-canonically''.
\end{remark}

\begin{definition}[The gluing $(2,1)$-category $\GlCoHA(X)$]
\label{def:glue-2cat}
\ClaimStatusDefinition
The \emph{gluing $(2,1)$-category} $\GlCoHA(X)$ has:
\begin{itemize}
\item \emph{Objects.} Pairs $(\mathfrak{U}, \sH_\bullet)$ with
 $\mathfrak{U}$ an analytic open cover whose charts satisfy (VdB)
 or (DM) of Proposition~\ref{prop:chart-quiver-presentation}, and
 $\sH_\bullet$ a simplicial $(2,1)$-sheaf of CoHAs on
 $(X, \mathfrak{U})$ satisfying (S1)--(S4) of
 Definition~\ref{def:simplicial-sheaf-cohas}.
\item \emph{1-morphisms.}
 $(\mathfrak{U}, \sH_\bullet) \to (\mathfrak{U}', \sH'_\bullet)$
 is a common refinement $\mathfrak{V}$ plus a homotopy-coherent
 simplicial 1-morphism
 $\sH_\bullet|_{\mathfrak{V}} \to \sH'_\bullet|_{\mathfrak{V}}$ in
 the Joyal model structure on simplicial sets valued in
 $\mathrm{Alg}^{\mathrm{FM}}_{(2,1)}$ (Lurie HA~4.7.5.1).
\item \emph{2-morphisms.} Homotopies of such homotopy-coherent
 1-morphisms on a further common refinement.
\end{itemize}
Composition is common-refinement composition at 1-morphism level
and vertical pasting at 2-morphism level. The poset of analytic
open covers under refinement is filtered, so common refinements
always exist and composition is well-defined up to contractible
choice.
\end{definition}

\begin{theorem}[Global CoHA as homotopy totalisation, conditional on stratum membership]
\label{thm:global-coha-limit}
\ClaimStatusProposition
Let $X$ be a Calabi--Yau threefold.
\begin{enumerate}
\item \emph{(Conditional non-emptiness.)} If $X$ belongs to at
 least one of the four equivariance strata of
 Section~\ref{sec:four-stratum-taxonomy}; toric, reduced
 $\bC^\times \times \Aut$, orbifold inertia, or lattice-polarised
 period; then $\GlCoHA(X)$ is non-empty.
\item \emph{(Torsor-pointed uniqueness.)} Every object
 of $\GlCoHA(X)$ has the same homotopy totalisation in
 $\mathrm{Alg}^{\mathrm{FM}}_{(2,1)}$:
 \[
 \sH(X) \;=\; \mathrm{holim}_{[n] \in \Delta}\, \sH_n,
 \]
 well-defined after fixing the coherent Keller--Yang/Tamarkin
 2-isomorphism torsor, by
 Proposition~\ref{prop:descent-reduction}.
\item \emph{(Étale functoriality.)} $X \mapsto \sH(X)$ extends to
 a functor from the $(2,1)$-category of Calabi--Yau threefolds
 and étale Calabi--Yau morphisms into
 $\mathrm{Alg}^{\mathrm{FM}}_{(2,1)}$.
\end{enumerate}
\end{theorem}

\begin{proof}[Proof of Theorem~\ref{thm:global-coha-limit}]
\emph{(1).} Each stratum supplies a geometric source for the
2-simplex filler $\eta_{\alpha\beta\gamma}$ of (S3):
\begin{itemize}
\item \emph{Stratum (i), toric.} On triple overlaps
 $U_{\alpha\beta\gamma}$ of toric charts, the DWZ mutation kernels
 $K_{\alpha\beta}, K_{\beta\gamma}, K_{\alpha\gamma}$ restrict to
 the same $T^d$-equivariant kernel on $U_{\alpha\beta\gamma}$, so
 $\eta_{\alpha\beta\gamma}$ is the canonical identity
 (Davison--Meinhardt arXiv:1601.02479, Theorem~A;
 Rapcak--Soibelman--Yang--Zhao arXiv:2007.13365, Theorem~1.5).
\item \emph{Stratum (ii), reduced $\Aut$.} $\Aut(X)$-equivariance
 promotes the mutation kernel to an $\Aut(X)$-equivariant kernel;
 equivariance plus Keller--Yang contractibility supplies
 $\eta_{\alpha\beta\gamma}$.
\item \emph{Stratum (iii), orbifold $X = [Y/G]$.} BKR
 arXiv:math/9908027 Theorem~1.1 furnishes the filler by McKay
 equivalence on the inertia stack $I(X)$.
\item \emph{Stratum (iv), lattice-polarised.} On the Humbert locus
 of the period domain, the Borcherds lift of the relevant Jacobi
 form (Borcherds arXiv:alg-geom/9609022) supplies an automorphic
 cocycle filling $\eta_{\alpha\beta\gamma}$.
\end{itemize}
Any single stratum-membership gives non-emptiness.
\emph{(2)} is Proposition~\ref{prop:descent-reduction}.
\emph{(3)} follows from base change for vanishing cycles
(Kashiwara, Publ.~RIMS 20 (1984), Chapter~5) and the étale
compatibility of the Hall product
(Davison arXiv:1311.7172, Section~5).
\end{proof}

\begin{remark}[Non-emptiness outside the four strata is open]
\label{rem:non-emptiness-conditional}
For compact Calabi--Yau threefolds with no torus action, no
residual automorphism, no orbifold presentation, and no
period-domain lift; the quintic, Schoen's bicubic, generic
complete intersections; the 2-simplex filler of (S3) has no
geometric source of any of the four types above; the
Davison--Meinhardt dimensional-reduction CoHA supplies 1-simplex
data but not coherent 2-simplex data for a global cover.
Non-emptiness of $\GlCoHA(X)$ for such $X$ is the open problem
formulated in \S\ref{sec:gluing:obstructions}.
\end{remark}

\begin{remark}[Relation to the categorical viewpoint]
\label{rem:cat-viewpoint-coha}
The object $\sH(X)$ produced by
Theorem~\ref{thm:global-coha-limit} is the \emph{chart-wise} or
\emph{combinatorial} global critical CoHA. It is the chain-level
incarnation of the $(\infty, 1)$-categorical CoHA of the
Kontsevich--Soibelman stack of objects in the Calabi--Yau
$\infty$-category $\Perf(X)$ equipped with its Serre-duality
cyclic structure. The $(\infty, 1)$-categorical and chain-level
constructions are equally load-bearing; neither subsumes the other, and
the chart-wise description exposed here is the one on which
explicit formulas live.
\end{remark}

\subsection{The four-stratum taxonomy}
\label{sec:four-stratum-taxonomy}

The critical question left open by
Theorem~\ref{thm:global-coha-limit} is whether $\GlCoHA(X)$ is
non-empty for a given $X$. Answering it requires naming the
\emph{equivariance stratum} on which $X$ sits: the symmetry group
whose cocycles assemble the overlaps. Four strata exhaust the
known mechanisms.

\paragraph{The four strata from the automorphism tower.} Fix
$(X, \Omega_X)$. The symmetry data available for a $\GlCoHA(X)$
gluing cocycle is exactly the \emph{tower of automorphism groups
that preserve $\Omega_X$ and act coherently on chart presentations}:
\[
 \mathrm{Aut}^{0}(X, \Omega_X)
 \;\lhd\;
 \mathrm{Aut}(X, \Omega_X)
 \;\lhd\;
 \mathrm{Aut}^{\widetilde{\phantom{X}}}(X, \Omega_X)
 \;\lhd\;
 \Gamma_{\mathrm{arith}}\bigl(\mathrm{Period}(X)\bigr).
\]
Each successive level adds a new source of gluing cocycle data
(continuous $\to$ discrete $\to$ orbifold-cover discrete $\to$
arithmetic). The four strata correspond one-to-one to the four
rows of this filtration.

\emph{(i) Toric.} $\mathrm{Aut}^0(X, \Omega_X) \supseteq T^d =
(\bC^\times)^d$: the identity component is a
volume-preserving torus. Examples: $\bC^3$, resolved conifold,
local $\bP^2$, banana threefold, SPP.

\emph{(ii) Reduced torus times automorphism.} $\mathrm{Aut}^0(X,
\Omega_X) \supseteq \bC^\times$ (the $\Omega_X$-scaling) and
$\pi_0 \mathrm{Aut}(X, \Omega_X)$ is a non-trivial discrete group
$G$ acting on $X$ itself. Examples: K3, K3~$\times$~E, abelian
surfaces, generic elliptic fibrations.

\emph{(iii) Orbifold inertia.} $X = [Y/G]$ with
$G \subset SU(d)$ finite acting on a smooth CY~$Y$; the symmetry
group $G$ does not act on $X$ directly, only on the cover~$Y$.
The inertia stack $I(X) = \coprod_{[g] \in G/\text{conj}}
Y^g / C_G(g)$ records the action of $G$ on $Y$ after the quotient.
Examples: $[\bC^3/\bZ_n]$, $[\bC^3/G]$ with $G \subset SU(3)$
finite, Mathieu $M_{24}$ acting on K3.

\emph{(iv) Lattice-polarised period.} The symmetry group acts
not on $X$ or on a cover of $X$, but on the period domain
$\mathcal{D}_L$ (for a fixed polarisation lattice
$L \subset H^2(X, \bZ)$); the group is the arithmetic subgroup
$\Gamma \subset O(L \otimes \bR)$ preserving $L$. Examples: K3 at
Humbert divisors, K3~$\times$~E via Gritsenko $\Delta_5$, Igusa
$\Phi_{10}$ for stable $K3^{[n]}$.

\begin{proposition}[The four strata exhaust the tower]
\label{prop:four-strata-exhaust}
\ClaimStatusProposition
The four strata above correspond to the four possible locations
along the tower at which the symmetry first appears:
(i) continuous identity component; (ii) discrete component acting
on $X$ directly; (iii) discrete component acting only on a cover
of $X$; (iv) arithmetic component acting on the period domain.
For a Calabi--Yau threefold $X$, the set of strata containing $X$
is determined by which levels of this tower are non-trivial; any
fifth stratum would correspond to a fifth level of the tower
outside $\mathrm{Aut}^0 \subset \mathrm{Aut} \subset
\mathrm{Aut}^{\widetilde{\phantom{X}}} \subset
\Gamma_{\mathrm{arith}}$, which exhausts the symmetry
classification of $(X, \Omega_X)$ (Deligne, \emph{Équations
différentielles à points singuliers réguliers}, Lecture Notes in
Math.~163 (1970), Chapter~II on monodromy for arithmetic rows
on period domains; Dolgachev arXiv:alg-geom/9502005 on
lattice-polarised K3 period maps).
\end{proposition}

\begin{proof}[Argument]
A symmetry producing 2-simplex-filler data for (S3) of
Definition~\ref{def:simplicial-sheaf-cohas} acts either (a) on $X$
itself preserving $\Omega_X$, or (b) on a cover $Y \to X$
preserving $\Omega_Y$, or (c) on the period domain of
$H^d(X, \bC)$. Case (a) partitions by connectedness of the acting
group: the identity component gives (i) when non-trivial, and
otherwise the non-trivial discrete component gives (ii). Case (b)
is (iii). Case (c) is (iv): Deligne, \emph{\'Equations
diff\'erentielles \`a points singuliers r\'eguliers}, LNM~163
(1970), Chapter~II, establishes that period-domain automorphisms
preserving an integral polarisation are arithmetic, so no
non-arithmetic symmetry acts on the period datum.
Dolgachev arXiv:alg-geom/9502005 Theorem~3.1 confirms this for
lattice-polarised K3. No further case remains: a hypothetical
symmetry of $(X, \Omega_X)$ lying outside (a)--(c) would fail to
act on any of the three categorical data on which (S3) is defined
(local charts, finite covers, period domain), and hence would not
define a 2-simplex filler. The tower is complete.
\end{proof}

\begin{theorem}[Four-stratum structure of gluing]
\label{thm:four-stratum-structural}
\ClaimStatusProposition
For each geometry $X$ to which the chart-wise construction of
Theorem~\ref{thm:global-coha-limit} applies, the gluing cocycle of
$\sH(X)$ factors through the strata (i)--(iv) of
Proposition~\ref{prop:four-strata-exhaust}. The theorem has three
parts: (A) stratum support, (C) cell enumeration; both
theorem-level; and (B) cocycle target, which is theorem-level
pointwise per stratum and conjectural at the uniform-across-strata
level specified in Remark~\ref{rem:bii-biv-scope}.
\begin{enumerate}
\item[(A)] \emph{Stratum support.} The set
\[
 \Str(X) \;:=\; \bigl\{ s \in \{(\text{i}), (\text{ii}),
 (\text{iii}), (\text{iv})\} :
 X \text{ carries a non-trivial cocycle in stratum } s \bigr\}
\]
is non-empty for every $X$ equipped with the $\bC^\times$-scaling
of $\Omega_X$ required for the Joyce--Song sign
$\varepsilon(E, F) = (-1)^{\chi(E, F)}$
(Joyce--Song arXiv:0810.5645, Definition~5.2 and Theorem~5.14;
Kontsevich--Soibelman arXiv:0811.2435, \S 6.3).
\item[(B)] \emph{Cocycle target, by stratum.} For $X$ of type
 $S$, the 2-simplex filler $\eta_{\alpha\beta\gamma}$ of (S3) in
 Definition~\ref{def:simplicial-sheaf-cohas} is witnessed by a
 first-cohomology class on the classifying datum associated with
 each stratum of $S$:
\begin{itemize}
\item[(i)] $H^1(X, \underline{\Pic^{T^d}})$, the torus-equivariant
 Picard sheaf, for toric $X$. (Theorem-level:
 Davison--Meinhardt arXiv:1601.02479, Theorem~A for the tilting
 cocycle; Rapcak--Soibelman--Yang--Zhao arXiv:2007.13365,
 Theorem~1.5 for its identification with the $T^d$-equivariant
 Picard class on triple toric overlaps.)
\item[(ii)] $H^1(X / \Aut_s(X), \underline{\Aut_s(X)})$, the
 $\Aut_s$-character cocycle, for $X$ with residual
 volume-preserving automorphism. (Theorem-level in chart-wise
 cases, e.g.\ K3 with Nikulin involution, via Mukai
 arXiv:math/0408129 Theorem~0.1; the full $K3 \times E$
 $\Aut_s$-cocycle is conjectural.)
\item[(iii)] $H^1(I(X/G), \underline{G})$, the inertia-stack
 Galois/McKay cocycle. (Theorem-level: BKR arXiv:math/9908027
 Theorem~1.1 plus Rapcak--Soibelman--Yang--Zhao arXiv:2007.13365
 Theorem~A.)
\item[(iv)] $H^1(\mathcal{D}_L / \Gamma_{\mathrm{arith}},
 \underline{\Gamma_{\mathrm{arith}}})$, the arithmetic automorphic
 cocycle on the period domain, equivalent to a Borcherds-lifted
 cusp form class. (Theorem-level case-wise: Borcherds
 arXiv:alg-geom/9609022 Theorem~10.1 for weak Jacobi lifts;
 Gritsenko--Nikulin arXiv:alg-geom/9504006 Theorem~2.1 for Siegel
 lifts; 8-form Gritsenko--Cl\'ery catalogue.
 A single unified $\Gamma_{\mathrm{arith}}$-valued statement
 across all such automorphic forms is conjectural.)
\end{itemize}
\item[(C)] \emph{Fifteen non-empty cells.} Of the $2^4 = 16$
 subsets $S \subset \{(\text{i}), (\text{ii}), (\text{iii}),
 (\text{iv})\}$, the subset $S = \varnothing$ is excluded by (A);
 the remaining 15 are the non-empty cells. Ten cells carry
 explicit witnesses (Example~\ref{ex:stratum-occupancy} below);
 five remain open for lack of a constructed geometry. For $X$ of
 type $S$ with $|S| \geq 2$, the cocycles indexed by distinct
 strata in $S$ are compatible via the overlap-coherence
 condition of
 Proposition~\ref{prop:overlap-coherence} below.
\end{enumerate}
\end{theorem}

\begin{proof}[Proof]
\emph{(A).} Suppose $X$ carries $\Omega_X$ nowhere-vanishing and
admits a chart-wise presentation with $\bC^\times$-scaling. The
critical CoHA construction requires a $\bC^\times$ action on
$\Omega_X$ to define the Joyce--Song sign
$\varepsilon(E, F) = (-1)^{\chi(E, F)}$
(Joyce--Song arXiv:0810.5645, Definition~5.2 introduces
$\varepsilon$; Theorem~5.14 proves the integration map is an algebra
homomorphism via the $\bC^\times$-compatibility of $\varepsilon$;
Kontsevich--Soibelman arXiv:0811.2435, \S 6.3: critical cohomology
uses the $\bC^\times$-weight on the vanishing-cycle sheaf); this
$\bC^\times$ is contained in $\Aut^0(X, \Omega_X)$ at the base of
the tower of Proposition~\ref{prop:four-strata-exhaust}. Either
$\Aut^0 \supseteq T^d$ (stratum (i)), or $\Aut^0 = \bC^\times$ and
at least one of $\pi_0 \Aut(X, \Omega_X),\;
\mathrm{Aut}^{\widetilde{\phantom{X}}}(X, \Omega_X),\;
\Gamma_{\mathrm{arith}}$ is non-trivial (one or more of strata
(ii)--(iv)). In every case $\Str(X) \neq \varnothing$.

\emph{(B).} Each entry is a direct citation. (i) Davison--Meinhardt
arXiv:1601.02479 Theorem~A identifies the tilting cocycle of a
toric CY with a class in the torus-equivariant Picard group.
(ii) Mukai arXiv:math/0408129 Theorem~0.1 identifies symplectic
$\Aut_s$-cocycles on K3 with a class in
$H^1(K3 / \Aut_s, \underline{\Aut_s})$. (iii) BKR plus RSYZ
identify the orbifold cocycle on $[Y/G]$ with a class in
$H^1(I(X/G), \underline{G})$. (iv) Borcherds and
Gritsenko--Nikulin identify the automorphic cocycle with a class
in $H^1(\mathcal{D}_L / \Gamma_{\mathrm{arith}},
\underline{\Gamma_{\mathrm{arith}}})$, realised by a
Borcherds-lifted Siegel or Jacobi form.

\emph{(C).} The cell count is $2^4 - 1 = 15$ after removing
$S = \varnothing$, which is excluded by (A).
Example~\ref{ex:stratum-occupancy} lists the 10 known witnesses;
the remaining 5 are open as stated, pending construction of a
geometry $X$ with the appropriate $\Str(X) = S$. Overlap
compatibility for $|S| \geq 2$ is
Proposition~\ref{prop:overlap-coherence}.
\end{proof}

\begin{remark}[Part (B)(ii), (B)(iv) pointwise-theorem vs uniform-conjectural]
\label{rem:bii-biv-scope}
Parts (B)(ii) and (B)(iv) are theorem-level for specific geometries:
(ii) for K3 with Nikulin involution, abelian surfaces, and local
K3-fibrations; (iv) for each Borcherds-lifted form
$\Delta_5, \Phi_{10}, \Delta_1, \ldots$ in the Gritsenko--Cl\'ery
8-form catalogue individually. A \emph{uniform} statement across
all residual $\Aut_s$ for (ii) and all Borcherds lifts for (iv) is
the residual unsolved content; not because individual cases
fail, but because the cocycle sheaves $\underline{\Aut_s(X)}$ and
$\underline{\Gamma_{\mathrm{arith}}}$ are not yet shown to form a
stack presentation unifying the cases. Parts (A) and (C) are
theorem-level outright; parts (B)(i) and (B)(iii) are theorem-level
via direct citation.
\end{remark}

\begin{remark}[Scope of Part (A); the general-compact case]
\label{rem:part-a-scope}
Part (A) of Theorem~\ref{thm:four-stratum-structural} asserts
$\Str(X) \neq \varnothing$ \emph{for every $X$ equipped with a
$\bC^\times$-scaling}. A hypothetical compact Calabi--Yau $X$ with
no continuous or discrete automorphism (generic quintic $\bP^4[5]$,
Schoen bicubic $\bP^4[3, 3]$) satisfies the $\bC^\times$-scaling
only via the weighted-projective dilation inherited from an
ambient projective space, and even then carries only the
Davison--Meinhardt fall-back; the resulting cocycle (if any) lives
in a degenerate stratum not computable by the four-stratum method.
This is not a counterexample to (A); it is the stratum-free
fall-back addressed by Remark~\ref{rem:fifteen-cells} and treated
in \S\ref{sec:gluing:obstructions}.
\end{remark}

\begin{proposition}[Overlap coherence as fibre-product cohomology]
\label{prop:overlap-coherence}
\ClaimStatusProposition
For $X$ with $\Str(X) = S$ and $|S| \geq 2$, the global gluing
cocycle of Theorem~\ref{thm:four-stratum-structural}(B) is a
class in the $(2,1)$-categorical fibre product
\[
 \kappa_X \;\in\;
 \lim_{s \in S} H^1\!\left(\,
 \mathrm{site}_s(X),\,
 \underline{\mathcal{F}_s(X)}\,
 \right)
\]
where the limit is the $(2,1)$-categorical limit of cohomology
groups of sheaves over sub-sites of $X$, each term is the cocycle
target of Theorem~\ref{thm:four-stratum-structural}(B) for its
stratum $s$, and $\underline{\mathcal{F}_s(X)}$ denotes the
corresponding classifying sheaf
($\underline{\Pic^{T^d}},\;
 \underline{\Aut_s(X)},\;
 \underline{G},\;
 \underline{\Gamma_{\mathrm{arith}}}$).
For $s, s' \in S$ two overlapping strata, the face-matching
datum between them is a compatible edge of the fibre-product
diagram; for $|S| \geq 3$, the triple-overlap compatibility is a
coherence 2-cell in the limit diagram. Higher coherences (3-cell
and above) are contractible by
Proposition~\ref{prop:glcoha-truncation-forced}.
\end{proposition}

\begin{proof}[Argument]
This is the $(2,1)$-categorical Čech overlap compatibility for a
cover of a site by the four sub-sites
$\{\mathrm{site}_s(X)\}_{s \in \{(\text{i}), \ldots, (\text{iv})\}}$:
the global cocycle is the 1-cell in the $(2,1)$-limit, and
compatibility on double/triple/quadruple overlaps is recorded as
2-cell data in the limit. Because
Proposition~\ref{prop:glcoha-truncation-forced} forces the ambient
category at $(2,1)$, all higher coherences are contractible and
the limit is computed as a fibre product in the
$(2,1)$-categorical sense. Direct unwinding exhibits the
diagrammatic structure, with the four face-matching edges
corresponding to the four pairwise intersections of strata and
the triangular 2-cells corresponding to the triple intersections.
At $|S| = 4$ (the master four-fold cell), the 2-cell data is a
tetrahedral coherence in the fibre-product cohomology; at
$|S| = 3$ it is a triangular coherence; at $|S| = 2$ it reduces
to a single edge.
\end{proof}

\begin{remark}[Strata overlap; the taxonomy is $\mathcal{P}(\{\text{i,ii,iii,iv}\})$, not a partition]
\label{rem:strata-not-partition}
The four strata are not a partition of the CY\textsubscript{3}
landscape: a single $X$ frequently sits in several strata, one per
non-trivial level of the automorphism tower
(Proposition~\ref{prop:four-strata-exhaust}). Local $\bP^2$ occupies
$\{(\text i), (\text{iii})\}$ (toric and $\bZ_3$-orbifold of
$\bC^3$ by BKR); K3 occupies $\{(\text{ii}), (\text{iii}),
(\text{iv})\}$ (reduced-torus automorphism, Mathieu $M_{24}$,
Kuga--Satake period); K3 $\times$ E occupies the same master triple
cell. The taxonomy is therefore the full power set
$\mathcal{P}(\{(\text{i}), (\text{ii}), (\text{iii}),
(\text{iv})\})$, with 16 cells
(Remark~\ref{rem:fifteen-cells}).
\end{remark}

\begin{remark}[The 16 cells of $\mathcal{P}(\{\text{i},\text{ii},\text{iii},\text{iv}\})$]
\label{rem:fifteen-cells}
The subset lattice $\mathcal{P}(\{(\text i), (\text{ii}),
(\text{iii}), (\text{iv})\})$ has $2^4 = 16$ cells. Define: a
CY\textsubscript{3} $X$ is of \emph{type $S$} if the strata
containing $X$ are exactly~$S$. The 16 cells split into:
\begin{itemize}
\item The empty cell $S = \emptyset$: $X$ has no level of the
 automorphism tower (Proposition~\ref{prop:four-strata-exhaust})
 non-trivial. Such $X$ are the \emph{stratum-free} compact
 CY\textsubscript{3}s; the chart-wise $\GlCoHA$-gluing
 construction does not apply
 (Remark~\ref{rem:non-emptiness-conditional}). The
 Davison--Meinhardt dimensional-reduction CoHA
 (Davison arXiv:1311.7172, Theorem~A) supplies 1-simplex data
 alone, without 2-simplex coherence, so it furnishes an algebra
 but not an object of $\GlCoHA(X)$.
\item The 15 non-empty cells $S \neq \emptyset$: at least one
 stratum is available, so Theorem~\ref{thm:global-coha-limit}(1)
 applies and $\GlCoHA(X)$ is non-empty.
\end{itemize}
The partition into 16 cells is exhaustive (power set of 4 strata);
the empty cell is the obstruction locus addressed in \S\ref{sec:gluing:obstructions}.
\end{remark}

\begin{example}[Per-cell realisability of the 16-cell taxonomy]
\label{ex:stratum-occupancy}
Each cell $S \in \mathcal{P}(\{(\text i), (\text{ii}), (\text{iii}),
(\text{iv})\})$ either carries a known realising geometry or is
currently open. The listing follows the lattice order:
\[
\begin{array}{l|l|l}
 S & \text{Realising geometry / status} &
 \text{Source} \\\hline
 \emptyset & \text{quintic } \bP^4[5],\ \text{Schoen bicubic} &
 \text{stratum-free fall-back only} \\
 \{\text{i}\} & \bC^3,\ \text{resolved conifold}, &
 \text{toric, non-compact} \\
 & \quad\text{banana, SPP} & \\
 \{\text{ii}\} & \text{abelian surface, }E\text{ elliptic} &
 \text{reduced-}\Aut\text{ only} \\
 \{\text{iii}\} & [\bC^3/G],\; G \subset SU(3) &
 \text{orbifold only} \\
 \{\text{iv}\} & \text{generic lattice-polarised K3} &
 \text{lattice, no extra sym} \\
 \{\text{i,ii}\} & \text{(open: toric + residual }\Aut) &
 \text{no known example} \\
 \{\text{i,iii}\} & \text{local } \bP^2 = \Tot(K_{\bP^2}) &
 \text{BKR: local }\bP^2 = [\bC^3/\bZ_3] \\
 \{\text{i,iv}\} & \text{(open)} & \text{no known example} \\
 \{\text{ii,iii}\} & \text{K3 with Nikulin involution} &
 \text{auto + quotient} \\
 \{\text{ii,iv}\} & \text{abelian surface with CM lattice} &
 \text{CM-polarised} \\
 \{\text{iii,iv}\} & \text{K3 with Mathieu } M_{24} &
 \text{orbifold + lattice} \\
 \{\text{i,ii,iii}\} & \text{(open)} & \text{no known example} \\
 \{\text{i,ii,iv}\} & \text{(open)} & \text{no known example} \\
 \{\text{i,iii,iv}\} & \text{(open)} & \text{no known example} \\
 \{\text{ii,iii,iv}\} & K3,\ K3 \times E & \text{master triple} \\
 \{\text{i,ii,iii,iv}\} & \text{(open: four-fold coincidence)} &
 \text{no known example}
\end{array}
\]
Of the 15 non-empty cells, 10 carry known realising
CY\textsubscript{3}s; 5 are open for lack of a constructed
example. The empty cell is the stratum-free locus where
$\GlCoHA$ does not apply.
\end{example}

\begin{remark}[K3 and K3 $\times$ E at the master triple cell]
\label{rem:k3e-master-case}
K3 and K3 $\times$ E both occupy the cell
$\{(\text{ii}), (\text{iii}), (\text{iv})\}$, with cocycles in
three of the four strata independently. On K3 $\times$ E all
three cocycles are non-degenerate: stratum (ii) supplies the
$\bC^\times$-scaling plus
$\Aut(K3) \times \Aut(E)$ data, stratum (iii) supplies the
Mathieu twining when the K3 factor carries an $M_{24}$-class
action, and stratum (iv) supplies the Borcherds lift of the
Gritsenko form $\Delta_5$ with $\kBKM(\mathfrak{g}_{\Delta_5}) = 5$.
The four $\kappa_\bullet$-invariants
$\kcat(K3 \times E) = 0$ (Künneth total space),
$\kch^{\mathrm{Heis}} = 3$ (chiral Heisenberg--Mukai),
$\kBKM(\mathfrak{g}_{\Delta_5}) = 5$ (Borcherds),
$\kfib = 24$ (Mukai-lattice rank of the K3 fibre) come from four
distinct constructions, with the toric stratum (i) entering
degenerately via the $\bC^3$ stalks at nodal points of the K3
fibre. \S\ref{sec:gluing:k3e} assembles the cocycles explicitly.
\end{remark}

\subsection{Abstract framework in one paragraph}
\label{sec:gluing:abstract:summary}

The global critical CoHA $\sH(X)$ is the $(2,1)$-categorical limit
of a simplicial $(2,1)$-sheaf of local critical CoHAs over the
nerve of any quiver-with-potential cover, with overlap data given
by Fourier--Mukai mutation cocycles. The existence of such a
cover and the consistency of the overlap cocycles is controlled
stratum by stratum, along the four equivariance classes
\emph{(i) toric, (ii) reduced torus $\times$ automorphism, (iii)
orbifold inertia, (iv) lattice-polarised period}. The taxonomy of
$(X, \Omega_X)$-automorphism-tower locations gives 16 cells in
$\mathcal{P}(\{(\text{i}), (\text{ii}), (\text{iii}), (\text{iv})\})$;
the 15 non-empty cells have $\GlCoHA(X)$ non-empty
(Theorem~\ref{thm:global-coha-limit}(1)), and the distinguished
empty cell is stratum-free. K3 and $K3 \times E$ share the master
triple cell $\{(\text{ii}), (\text{iii}), (\text{iv})\}$;
$K3 \times E$ carries four mutually-compatible gluing cocycles,
assembled in \S\ref{sec:gluing:k3e}. Stratum-free compact CY\textsubscript{3}s
(the quintic, Schoen's bicubic) admit only the Davison--Meinhardt
fall-back; the obstruction to chart-wise assembly for them is the
subject of \S\ref{sec:gluing:obstructions}.

\S\ref{sec:gluing:toric-cech} realises the abstract framework chart-wise on stratum (i),
beginning with the two-chart conifold.
